\documentclass[11pt]{article}

\usepackage[margin=1in]{geometry}
\usepackage{amsmath, amssymb}
\usepackage{booktabs}
\usepackage{graphicx}
\usepackage{hyperref}
\usepackage{enumitem}

\title{Learning Flow Maps of Jump-Diffusion SDEs:\\
A Runaway-Electron Target and a 1D Verification}
\author{}
\date{\today}

\begin{document}
\maketitle

\noindent\textbf{Overview.} The eventual target is a 2D runaway-electron (RE)
model whose dynamics are naturally a stochastic differential equation (SDE) with
jumps on a bounded momentum domain. We plan to learn its large-time stochastic
flow map with a training-free diffusion model, plus an exit indicator for the
boundary. Before the physics problem we verify the approach on a 1D bounded
jump-diffusion, comparing the learned model against the classical numerical
method in accuracy, speed, and long-horizon trajectory prediction.

%======================================================================
\section{The eventual target: a 2D runaway-electron model}
%======================================================================

\paragraph{The model.} We track a single relativistic electron in 2D momentum
space,
\begin{equation}
X_t = (p,\ \xi), \qquad p=\text{normalized momentum}, \quad
\xi=\cos\theta=\text{pitch-angle cosine},
\end{equation}
driven by the toroidal electric field, collisional drag and scattering, and rare
large-angle collisions. The dynamics take the jump-diffusion form
\begin{equation}
\label{eq:jdsde}
dX_t = \underbrace{b(X_t)\,dt}_{\text{drift}}
     + \underbrace{\sigma(X_t)\,dW_t}_{\text{diffusion}}
     + \underbrace{\int \gamma(X_{t^-},z)\,\tilde N(dt,dz)}_{\text{jumps}},
\end{equation}
with the three pieces given concretely by
\begin{align}
b_p &= \hat E\,\xi - F_{\rm drag}(p), &
F_{\rm drag}(p) &\sim \tfrac{1+p^2}{p^2}, \\
b_\xi &= \frac{\hat E\,(1-\xi^2)}{p}, &
d\xi &\supset \sqrt{2\,\nu_D(p)\,(1-\xi^2)}\;dW_\xi ,
\end{align}
where $\hat E=E/E_c$ is the normalized field, $F_{\rm drag}$ is the collisional
friction (it falls off as $\sim 1/v^2$ at high $p$, which is why runaways exist),
and $\nu_D$ is the pitch-angle scattering rate. The momentum domain is bounded by
the \emph{runaway separatrix} $p_c$: below it the electron rejoins the thermal
bulk (an exit/absorbing boundary).

\paragraph{Why it is an SDE with jumps.} The full Coulomb collision operator
splits into a small-angle part and a large-angle part. The small-angle
(grazing) collisions accumulate into the Gaussian diffusion $\sigma\,dW$ (the
classical Fokker--Planck picture). The rare large-angle, close (``knock-on'')
collisions transfer a large momentum chunk in a single event; these are discrete,
Poisson-timed, and heavy-tailed (M{\o}ller cross-section
$d\Sigma/d\epsilon\propto\epsilon^{-2}$), i.e.\ exactly the jump term in
Eq.~\eqref{eq:jdsde}~\cite{embreus2018}. The same large-angle collisions drive
the runaway \emph{avalanche}~\cite{rosenbluth1997}; here we keep the
single-particle jump-diffusion and treat the avalanche as the motivating context.
Framing runaway generation through the separatrix $p_c$ as an exit problem
matches existing RE methodology~\cite{delcastillo_exit}.

\paragraph{What we plan to do.} Rather than integrating Eq.~\eqref{eq:jdsde} with
many small time steps, we will learn its \emph{stochastic flow map} --- the
one-step transition density $p(X_{t+\Delta t}\mid X_t)$ at a \emph{large}
$\Delta t$ --- with a training-free diffusion model~\cite{liu2024,yang2025}, plus
an exit indicator for the $p_c$ boundary~\cite{yang2025}. One evaluation of the
learned map then advances the state by a large step, replacing many integrator
sub-steps. The new ingredient relative to existing diffusion flow-map work is the
\emph{jump} term, which is the subject of the verification below.

%======================================================================
\section{A 1D verification example}
%======================================================================

We verify the approach on a 1D jump-diffusion in a \emph{bounded} domain with an
absorbing boundary --- the structure the 2D runaway-electron problem needs (jumps,
a large-$\Delta t$ flow map, and exit across the separatrix). We learn two
objects, a flow map for the in-domain transition and an \emph{exit indicator} for
the boundary, and compare against the classical numerical method in accuracy,
speed, and long-horizon trajectory prediction.

\paragraph{The model.} A 1D jump-diffusion on $[0,6]$ with absorbing ends,
\begin{equation}
dX_t = \sigma_b\,dW_t + dJ_t, \qquad X_t \in [0,6],
\end{equation}
with $J_t$ a compound-Poisson process (rate $1$, jumps $Y\sim\mathcal N(0,0.7^2)$),
$\sigma_b=0.7$, and a large step $\Delta t=1$. A particle that crosses either end
leaves the domain (exits). The ground truth is a fine-step jump-adapted
Euler--Maruyama integrator with the same boundaries.

\paragraph{The exit-indicator idea.} Surviving particles stay inside $[0,6]$ by
construction, so the flow map is only ever evaluated in-domain --- over
arbitrarily long rollouts. We learn the boundary with a binary classifier
$\mathcal E(x)=\Pr(\text{exit during }\Delta t \mid x)$, trained by cross-entropy
on labeled survive/exit data (the JCP exit indicator of~\cite{yang2025}). One
coarse step is \emph{exit, then propagate}: $\mathcal E$ removes the particles
that cross the boundary, and the flow map --- trained only on \emph{surviving}
transitions --- advances the rest. Over a long rollout this tracks both the
decaying surviving population and the surviving distribution, with the flow map
always in-domain. This mirrors the runaway-electron separatrix $p_c$, where
particles leave the runaway region into the thermal bulk.

\paragraph{What we did.} Generate large-$\Delta t$ data from the bounded SDE:
survivor pairs $(X_t,X_{t+\Delta t})$ and exit labels $(X_t,\text{exit?})$. Learn
the survivor flow map with the training-free method --- learn the increment
$X_{t+\Delta t}-X_t$ via a probability-flow ODE with a training-free Monte-Carlo
score, distilled into a small network --- and the exit indicator by cross-entropy.
At inference, roll the two heads forward and compare against the numerical method.

\paragraph{Accuracy (one step).} Figure~\ref{fig:b1} verifies the two heads
in-domain: the exit indicator recovers the U-shaped exit probability (high at the
walls, near zero mid-domain), and the survivor transition density matches the
numerical truth ($W_1=0.004$ and $0.015$ at $x=1.5,3.0$).

\begin{figure}[h]
\centering
\includegraphics[width=\textwidth]{jump1d/figs/B1_onestep.png}
\caption{One step. Left: exit probability $\Pr(\text{exit}\mid x)$, exit indicator
(red) vs.\ numerical (black). Middle/right: survivor transition density, flow map
(red) vs.\ numerical survivors (gray), at $x=1.5,3.0$.}
\label{fig:b1}
\end{figure}

\paragraph{Long-horizon trajectory and speed.} Figure~\ref{fig:b2} propagates to
$T=8$ ($K=8$ big ML steps vs.\ the integrator's $8\times200=1600$ sub-steps). The
ML rollout tracks the \emph{surviving fraction} at every step (final $0.316$
vs.\ $0.318$, a $0.6\%$ error) and reproduces the \emph{terminal surviving
density} ($W_1=0.010$, near the sampling floor), about $40\times$ faster than the
numerical method. The structural reason for the speedup is that the learned map
advances the state in one big step rather than many sub-steps.

\begin{figure}[h]
\centering
\includegraphics[width=\textwidth]{jump1d/figs/B2_rollout.png}
\caption{Long rollout to $T=8$. Left: surviving fraction vs.\ step, ML rollout
(red) vs.\ numerical (black). Right: terminal surviving density at $T=8$, ML (red)
vs.\ numerical (gray), $W_1=0.010$. The ML map takes $8$ big steps vs.\ the
integrator's $1600$ small steps. (Wall-clock is indicative: integrator on CPU,
network on GPU; the hardware-independent statement is $8$ vs.\ $1600$ steps.)}
\label{fig:b2}
\end{figure}





\newpage

%======================================================================
\section{Problem settings for independent reproduction}
\label{sec:settings}

This section is written to be handed out on its own: it is self-contained,
restates every equation it needs, and cites the literature in full. It
poses four problems on learning large-time stochastic flow maps of
jump-diffusion SDEs, in increasing order of difficulty --- each problem
adds one ingredient to the previous one. Only the problem settings are
given: the model with every parameter value, the reference numerical
method, the objects to learn, the tests to run, and the results to expect
(with the error magnitudes we measured in our own implementation). No
implementation is provided --- build your own codes for the reference
integrator, the data generation, the learned model, and the comparisons;
at most, the training-data generation code may be shared as a starting
point. Two remarks. First, treat the quoted numbers as targets of
\emph{magnitude}, not answers: Monte-Carlo noise, training stochasticity,
and implementation choices all move them. Second, the method itself is
not a finished principle. Where your results deviate from the
expectations below, or where the construction feels fragile --- heavy
tails, steep exit fronts, the calibration of small exit probabilities ---
that is a starting point for independent research, not necessarily an
error to fix.

\paragraph{The task, in general.} All four problems concern a
jump-diffusion SDE
\begin{equation*}
dX_t = b(X_t)\,dt + \sigma(X_t)\,dW_t + dJ_t,
\end{equation*}
with $J_t$ a compound-Poisson jump process, on a domain $D$ that is
either bounded with absorbing boundaries (Problems 1--3) or periodic
(Problem 4). Instead of integrating with many small time steps, learn the
\emph{one-step transition law at a large coarse step} $\Delta t$ directly
from data: a generative map $G(x,z)$, $z\sim\mathcal N(0,I)$, such that
$G(x,z)$ is distributed like $X_{t+\Delta t}$ given $X_t=x$. On a bounded
domain the one-step kernel is a sub-probability kernel (particles can
exit), which factorizes into two learning problems,
\begin{equation*}
P_{\Delta t}(x,A)\;=\;\bigl(1-\mathcal E(x)\bigr)\,
P^{\rm surv}_{\Delta t}(x,A), \qquad A\subseteq D,
\end{equation*}
where $\mathcal E(x)=\Pr(\text{exit during }\Delta t\mid X_t=x)$ is a
classifier (the \emph{exit indicator}) trained on exit labels, and
$P^{\rm surv}_{\Delta t}$ is the flow map trained on \emph{surviving}
transitions only. One coarse step of the learned model is \emph{exit,
then propagate}: the indicator removes the exiting particles, the flow
map advances the rest. Since survivors stay in $D$ by construction, the
flow map is never evaluated out-of-domain, over arbitrarily long
rollouts. We use the training-free diffusion flow map of Y.~Liu, M.~Yang,
Z.~Zhang, F.~Bao, Y.~Cao, G.~Zhang, \emph{J.~Machine Learning for
Modeling and Computing} \textbf{5}(1), 25 (2024) and M.~Yang, Y.~Liu,
D.~del-Castillo-Negrete, Y.~Cao, G.~Zhang, \emph{J.~Comput.\ Phys.}
114434 (2025); any conditional generative model is admissible, and
comparing families is itself a worthwhile experiment.

\paragraph{Common setup.}
\begin{itemize}[leftmargin=1.4em]
\item \emph{Reference method (and data generator).} Jump-adapted
  Euler--Maruyama: split the coarse step $\Delta t$ into $n_{\rm sub}$
  sub-steps, draw the jump events of the compound-Poisson term by Poisson
  sampling (thinning where the rate is state-dependent), integrate drift
  and diffusion between jumps, and, on bounded domains, test absorption at
  every sub-step. Before judging any learned model, verify the
  \emph{reference's own} convergence by comparing at least two values of
  $n_{\rm sub}$ against a finer one on the observables you will report.
\item \emph{Training data.} Pairs $(X_t, X_{t+\Delta t})$ generated by
  the reference integrator from starts sampled over the domain, one
  coarse step each; on bounded domains, split into survivor pairs (for
  the flow map) and exit labels (for the indicator).
\item \emph{Evaluation.} Compare distributions by the Wasserstein-1
  distance of marginals (estimate the sampling floor by splitting the
  reference sample in half), exit and fate probabilities by their
  absolute errors against the Monte-Carlo reference, and always report
  the step-count ratio (coarse steps vs.\ integrator sub-steps) alongside
  any wall-clock speedup.
\end{itemize}

%----------------------------------------------------------------------
\subsection{Problem 1: 1D bounded jump-diffusion, constant coefficients}
%----------------------------------------------------------------------

\paragraph{Setting.}
\begin{equation*}
dX_t = \sigma_b\,dW_t + dJ_t \quad\text{on } [0,6],
\qquad \text{absorbing at both ends,}
\end{equation*}
with $\sigma_b=0.4$ and $J_t$ compound Poisson with rate $\lambda=1$ and
jump sizes $Y\sim\mathcal N(\mu_Y,\sigma_Y^2)=\mathcal N(1.5,\,0.3^2)$.
Coarse step $\Delta t=1$; horizon $T=4$ ($K=4$ coarse steps) --- about
the longest horizon at which both test starts below retain usable
survivor statistics. Reference: $n_{\rm sub}=200$ sub-steps per coarse
step; use $\ge 2\times10^5$ particles per test. Suggested training set:
$4\times10^5$ coarse-step transitions from uniformly sampled starts.

\paragraph{Tests.} (B0) The one-step transition from the fixed starts
$x_0=1$ and $x_0=5$, plotted as \emph{sub-probability} densities
(normalized by the particles started, so exit mass appears as missing
area); in free space (no boundary) the density has the closed form of a
jump-count mixture,
\begin{equation*}
p_{\Delta t}(y\mid x_0)=\sum_{n\ge0}\mathrm{Pois}(n;\lambda\Delta t)\;
\mathcal N\!\bigl(y;\ x_0+n\mu_Y,\ \sigma_b^2\Delta t+n\sigma_Y^2\bigr),
\end{equation*}
--- overlay it. (B1) The exit-probability curve
$\Pr(\text{exit}\mid x)$ over the domain and the survivor densities at
$x=1.5,\,3$. (B2) The rollout to $T=4$ from both starts: surviving
fraction at every coarse step and the terminal surviving density.

\paragraph{What to expect.} A strongly multi-modal one-step density with
modes near $x_0+1.5n$; an \emph{asymmetric} exit curve --- diffusive near
the lower wall, rising well before the upper wall, where exits are
jump-driven and nonlocal; survival at $T=4$ of $\approx0.35$ from $x_0=1$
and $\approx0.02$ from $x_0=5$. A well-trained flow map reaches survivor
$W_1\sim10^{-2}$. The hard part is the jump-driven exit layer near the
upper wall: our exit indicator gives $0.575$ vs.\ the true $0.615$ at
$x_0=5$, and that one-step underestimate compounds visibly over the
rollout (final survival $0.032$ vs.\ $0.023$). The step-count ratio is
$4$ vs.\ $800$; our wall-clock speedup was $\approx78\times$.

%----------------------------------------------------------------------
\subsection{Problem 2: 1D bounded jump-diffusion, state-dependent
coefficients}
%----------------------------------------------------------------------

\paragraph{Setting.} Same domain, boundaries, coarse step, learned
objects, reference, sample sizes, and tests (from the same starts
$x_0=1,5$), but every coefficient now depends on the state:
\begin{equation*}
dX_t = \kappa\,(x^*-X_t)\,dt + \sigma(X_t)\,dW_t + dJ_t^{\lambda(X_t)}
\quad\text{on } [0,6],
\end{equation*}
with $\kappa=0.5$, $x^*=2.5$, $\sigma(x)=0.3+0.05x$, jump rate
$\lambda(x)=1.3-0.15x$, and jump sizes
$Y\mid x\sim\mathcal N(1.5-0.08x,\,0.3^2)$. There is no closed-form
overlay here --- the fine-step histogram is the reference. The horizon is
$T=8$ ($K=8$): the mean-reverting drift keeps survival healthy, so the
compositional test can run twice as long as in Problem~1. This is the
first problem where the conditioning on $x$ is actually exercised (and
where the classical method has no large-step shortcut).

\paragraph{What to expect.} The one-step law varies visibly with the
state: bimodal from $x_0=1$ (high rate, large kicks), a single
drift-shifted bump from $x_0=5$ (low rate, downward drift). Survival at
$T=8$ of $\approx0.63$ from $x_0=1$ and $\approx0.45$ from $x_0=5$;
terminal $W_1$ near the sampling floor (a few $10^{-3}$). Note the
expectation: this example should come out \emph{more} accurate than
Problem~1, because the accuracy-limiting feature there is the sharp exit
layer, not the conditioning. If your state-dependent run is much worse
than your constant-coefficient one, inspect the conditioning (input
normalization, coverage of the training starts) before anything else.

%----------------------------------------------------------------------
\subsection{Problem 3: a 2D runaway-electron model}
%----------------------------------------------------------------------

\paragraph{Background.} A relativistic electron in a plasma with an
electric field above the critical field can ``run away'': collisional
friction falls off with velocity, so a fast-enough electron accelerates
without bound. Small-angle Coulomb collisions act as drift and diffusion;
rare large-angle (``knock-on'') collisions transfer a large momentum
fraction in single events --- exactly a jump term. The state is
$X=(p,\xi)$ with $p$ the momentum in units of $m_ec$ and
$\xi=\cos\theta$ the pitch-angle cosine; time is in relativistic
collision times $\tau_c$, the field $\hat E=E/E_c$ in units of the
critical field.

\paragraph{Setting: drift and diffusion.} The test-particle model of
G.~Zhang and D.~del-Castillo-Negrete, \emph{Phys.\ Plasmas} \textbf{24},
092511 (2017), Eqs.~(5)--(6):
\begin{align*}
dp &= \Bigl[\hat E\,\xi \;-\; \bar f\,\frac{1+p^2}{p^2} \;-\;
      \frac{\gamma p\,(1-\xi^2)}{\tau_{\rm syn}}\Bigr]\,dt \;+\; dJ_t,\\
d\xi &= \Bigl[\frac{\hat E\,(1-\xi^2)}{p} \;+\;
      \frac{\xi(1-\xi^2)}{\tau_{\rm syn}\gamma} \;-\; \xi\,\nu_c\Bigr]\,dt
      \;+\; \sqrt{\nu_c\,(1-\xi^2)}\;dW \;+\; dJ^\xi_t,
\qquad \nu_c=\bigl(Z+\bar f\,\bigr)\frac{\gamma}{p^3},
\end{align*}
with $\gamma=\sqrt{1+p^2}$; note $dp$ carries no Brownian term --- all
momentum-space randomness enters through the jumps. Parameters:
$\hat E=2$, $Z=1$, $\ln\Lambda=15$, $\tau_{\rm syn}=100$. The factor
$\bar f=\overline{\ln\Lambda}/\ln\Lambda$ is defined with the jump term
below.

\paragraph{Setting: the knock-on jump process.} All formulas from
O.~Embr\'eus, A.~Stahl, T.~F\"ul\"op, \emph{J.~Plasma Phys.} \textbf{84},
905840102 (2018):
\begin{itemize}[leftmargin=1.4em]
\item \emph{Rate}: $\lambda(p)=\tfrac12 nv\,\sigma(p)$, with $\sigma(p)$
  the M{\o}ller cross section integrated over energy transfers
  $\epsilon\in[\epsilon_{\min},(\gamma-1)/2]$ --- an analytic closed form,
  their Eq.~(35); the upper limit is the indistinguishability bound. In
  $\tau_c$ units: $\lambda\tau_c=R\,\beta\,\hat\sigma/(4\ln\Lambda)$ with
  $\hat\sigma=\sigma/(2\pi r_0^2)$, $\beta=p/\gamma$, and
  $R=n_{\rm tot}/n_e=1$ here.
\item \emph{Jump law}: the transferred energy $\epsilon$ follows the
  M{\o}ller spectrum, their Eq.~(5), heavy-tailed
  $\propto\epsilon^{-2}$. The retained electron (the more energetic
  outgoing branch, guaranteed by $\epsilon\le(\gamma-1)/2$) has
  $\gamma'=\gamma-\epsilon$ and is deflected by
  $\cos\theta_d=\sqrt{(\gamma'-1)(\gamma+1)/\bigl((\gamma'+1)(\gamma-1)\bigr)}$,
  their Eq.~(6), with uniform azimuth.
\item \emph{No double counting}: transfers below $\epsilon_{\min}$ remain
  in the drift-diffusion coefficients, whose electron-electron Coulomb
  logarithm is reduced,
  $\overline{\ln\Lambda}=\ln\Lambda-\tfrac12
  \ln[(\gamma-\gamma_m)/(\gamma_m-1)]$ with
  $\gamma_m=1+\epsilon_{\min}$, their Eqs.~(13)--(14). The reduction
  factor $\bar f$ multiplies the drag and the electron-electron unit of
  $\nu_c$ (not the electron-ion $Z$ part, which has no jump channel), as
  written in the equations above. Take $\epsilon_{\min}=0.02$.
\end{itemize}

\paragraph{Setting: domain and problem sizes.} $(p,\xi)\in[1,10]\times[-1,1]$,
absorbing at both $p$-ends ($p<1$: thermalization; $p>10$: confirmed
runaway), reflecting at $\xi=\pm1$ (the pitch diffusion vanishes there).
Coarse step $\Delta t=1\,\tau_c$; horizon $T=10\,\tau_c$ ($K=10$). The
exit indicator now has \emph{three classes}: survive / exit $p_{\min}$ /
exit $p_{\max}$. Suggested training set: $6\times10^5$ coarse-step
transitions from uniform starts.

\paragraph{Reference discipline.} This problem is where validating the
reference becomes serious. (i)~Scan $n_{\rm sub}$: we found
$n_{\rm sub}=200$ insufficient (a localized pitch-distribution bias at
low-$p$ starts, where scattering is fastest) and used $400$, verified
against $800$. (ii)~Scan the modeling split
$\epsilon_{\min}\in\{0.01,0.02,0.04\}$: expect genuine drifts confined to
the two absorbing-boundary layers --- the thermalization probability
rises with $\epsilon_{\min}$ and the one-step runaway probability falls,
both at the few-percent level, with drifts $<0.01$ elsewhere. Understand
why before quoting exit probabilities (hint: the
$\overline{\ln\Lambda}$ correction matches the \emph{mean} energy loss
across the split, but near an absorbing wall the exit probability is
sensitive to the fluctuation structure of that loss, not only its mean).

\paragraph{Tests.} (B0) One-step transition from $(p_0,\xi_0)=(2,0.6)$
(a contested start: acceleration, scattering, and jump losses compete)
and $(5,0.9)$ (runaway-bound), as 2D sub-probability densities with $p$-
and $\xi$-marginals. (B1) One-step exit maps over the domain (total exit
probability and the runaway channel separately), on a grid of roughly
$36\times18$ starts with $\gtrsim1500$ particles each. (B2) Rollout to
$T=10$: survival and the cumulative runaway/thermalization split per
coarse step, and the terminal surviving density.

\paragraph{What to expect.} Unlike the Gaussian-jump 1D problems, the
one-step density is \emph{uni-modal} --- a drift-shifted core with a
heavy downward tail in $p$; the jump signature sits in the tail and in
the exit mass, not in separated modes. The runaway exit map has a steep
S-shaped front in $(p,\xi)$. Fate split at $T=10$: from $(2,0.6)$
approximately survival $0.54$, runaway $0.33$, thermalized $0.14$; from
$(5,0.9)$ runaway $\approx0.98$. Survivor marginals at one step reach
$W_1\sim10^{-2}$ in $p$ and a few $10^{-3}$ in $\xi$. Expect the
\emph{exit indicator to be the accuracy-limiting component}, with the
largest errors on the steep runaway front, and expect naive uniform
training data to be insufficient there --- we needed exit labels
concentrated near the boundaries and wall-adapted input features; finding
your own remedy is encouraged. Once the model works, ensemble
computations become cheap: the runaway-probability map
$P_{\rm RE}(p_0,\xi_0)$ (roll a grid of starts to fate saturation,
$T=30\,\tau_c$), scans over initial seed distributions, and
first-passage-time distributions are natural extensions --- with the fine
integrator used only at a few anchor points.

%----------------------------------------------------------------------
\subsection{Problem 4: a 3D periodic transport benchmark}
%----------------------------------------------------------------------

\paragraph{Setting.} A passive tracer on the torus
$\mathbb T^3=[0,2\pi)^3$:
\begin{equation*}
dX_t = \mathrm{Pe}\,v(X_t)\,dt + dW_t + dJ_t,
\end{equation*}
with the Arnold--Beltrami--Childress velocity field
\begin{equation*}
v(x,y,z)=\bigl(A\sin z + C\cos y,\;\; B\sin x + A\cos z,\;\;
C\sin y + B\cos x\bigr), \qquad A=B=C=1,
\end{equation*}
at P\'eclet number $\mathrm{Pe}=5$; $W_t$ a standard 3D Brownian motion;
and $J_t$ compound Poisson with rate $\lambda=4$, jumps
$\Delta J=\ell\,\Omega$ with $\Omega$ uniform on the unit sphere and the
length $\ell$ drawn from the exponentially tempered, truncated power law
\begin{equation*}
f(\ell)\;\propto\;\ell^{-1-\alpha}\,e^{-\ell/\ell_c}
\quad\text{on }[\ell_{\min},\ell_{\max}],
\qquad \alpha=1.5,\;\ell_{\min}=0.2,\;\ell_c=1,\;\ell_{\max}=2.5
\end{equation*}
(tempering: A.~Cartea, D.~del-Castillo-Negrete, \emph{Phys.\ Rev.\ E}
\textbf{76}, 041105 (2007); truncation: R.~N.~Mantegna, H.~E.~Stanley,
\emph{Phys.\ Rev.\ Lett.} \textbf{73}, 2946 (1994); $\ell_{\max}<\pi$, so
a single jump cannot wrap the box). Coarse step $\Delta t=0.1$
($\lambda\Delta t=0.4$ jumps per step); horizon $T=3$ ($K=30$). The
domain is periodic --- no exit indicator; the learned object is the full
transition kernel. Reference: $n_{\rm sub}=200$, scanned against $100$
and $400$. Suggested training set: $3\times10^6$ one-step pairs from
uniform starts (the field is divergence-free, so the uniform density is
invariant and this covers the conditioning domain exactly).

\paragraph{Two companion models}, trained separately: the \emph{no-jump
twin} ($J\equiv0$) and the \emph{walk} variant --- rate
$\lambda(x)=4\,|v(x)|^2/\langle|v|^2\rangle$ (flights initiated in the
fast jet regions; $\langle|v|^2\rangle=A^2+B^2+C^2$ under the uniform
measure), direction $+v(x)/|v(x)|$, same length law but
$\ell_{\max}=1.5$. The tighter cutoff matters: flow-aligned flights
compound coherently with the advection within one coarse step, and with
$\ell_{\max}=2.5$ the one-step displacement becomes ambiguous on the
torus at a measurable rate. \emph{Measure} your mis-wrap fraction rather
than assuming it is zero.

\paragraph{A design point to solve, not copy.} The conditioning must
respect the periodicity --- naive conditioning on $x\in[0,2\pi)^3$
breaks at the seam. We condition on $(\cos x_i,\sin x_i)_{i=1}^3$ and
learn minimal-image increments
$\bigl((y-x+\pi)\bmod 2\pi\bigr)-\pi$; other constructions are possible
and worth comparing.

\paragraph{Tests.} From the fixed starts $(\pi/2,\pi/2,\pi/2)$ and
$(0,\pi,\pi/2)$, with displacements accumulated as sums of minimal-image
coarse-step increments: one-step increment distributions per axis on a
log scale (the tails carry the jumps); over the $30$ steps, the
eigenvalues of the displacement covariance divided by $2t$ (finite-time
dispersion --- check, do not assume, a plateau), the non-Gaussianity
parameter
$\mathrm{NGP}(t)=\tfrac{3\langle|D-\langle D\rangle|^4\rangle}
{5\langle|D-\langle D\rangle|^2\rangle^2}-1$ (zero for Gaussian
transport), a mixing diagnostic toward uniformity (e.g.\ the $L_1$
distance of the 2D marginals from uniform), preservation of the uniform
invariant density under the learned map, and the terminal displacement
distributions. Run all three models and compare.

\paragraph{What to expect.} The three-model comparison is the point of
the benchmark, and it has a sign structure worth predicting before
computing: the finite-time dispersion trace orders as \emph{isotropic
jumps $<$ no jumps $<$ flow-aligned flights} ($\approx36<38<48$ at the
symmetric start by $T=3$) --- isotropic relocations \emph{decorrelate}
particles from coherent ballistic advection, removing more dispersion
than the jump variance adds, while aligned flights \emph{extend} it.
Jumps make the one-step tails much heavier ($99.9\%$ displacement
quantile $\approx2.3$ vs.\ $1.6$--$1.8$ without jumps). Expect the
learned map's weakest point to be exactly this extreme tail: rare large
jumps are underrepresented in any finite training set, so the
$10^{-3}$--$10^{-4}$ tail will come out under-resolved. The uniform
density should be preserved at the histogram noise floor. A good learned
result reproduces the dispersion ordering above with each trace within a
few percent, and tracks $\mathrm{NGP}(t)$ after the first few steps.

\paragraph{Beyond reproduction.} Open directions that start from these
problems: exit-indicator calibration near steep fronts (rare-event
training, alternative losses); tail-aware training for heavy-tailed jump
laws; a quantitative consistency analysis of the learned kernel (when and
how fast does it converge to the true one as the data grow?); replacing
the training-free diffusion map by other conditional generative families;
adaptive coarse steps; and, on the physics side, the avalanche extension
of the runaway-electron model (the same jump machinery with particle
creation). Improvements on any of these are new results, not
reproductions.




\paragraph{Conclusion.} On a 1D bounded jump-diffusion, the learned flow map plus
exit indicator reproduce both the survival dynamics and the terminal distribution
over a long horizon, at a large $\Delta t$ and about $40\times$ faster than the
classical integrator. This exercises every ingredient the 2D runaway-electron
model of Eq.~\eqref{eq:jdsde} requires --- jumps, a large-$\Delta t$ flow map, and
an exit indicator for the separatrix --- before scaling up.

%======================================================================
\begin{thebibliography}{9}

\bibitem{liu2024} Y.~Liu, M.~Yang, Z.~Zhang, F.~Bao, Y.~Cao, G.~Zhang,
``Diffusion-model-assisted supervised learning of generative models for density
estimation,'' \textit{J. Machine Learning for Modeling and Computing}
\textbf{5}(1), 25--38 (2024). doi:10.1615/JMachLearnModelComput.2024051346.

\bibitem{yang2025} M.~Yang, Y.~Liu, D.~del-Castillo-Negrete, Y.~Cao, G.~Zhang,
``Generative AI models for learning flow maps of stochastic dynamical systems in
bounded domains,'' \textit{J. Comput. Phys.}, 114434 (2025).

\bibitem{rosenbluth1997} M.~N.~Rosenbluth, S.~V.~Putvinski,
``Theory for avalanche of runaway electrons in tokamaks,''
\textit{Nucl. Fusion} \textbf{37}, 1355 (1997).

\bibitem{embreus2018} O.~Embr\'eus, A.~Stahl, T.~F\"ul\"op,
``On the relativistic large-angle electron collision operator for runaway
avalanches in plasmas,'' \textit{J. Plasma Phys.} \textbf{84}, 905840102 (2018).

\bibitem{delcastillo_exit} M.~Yang, G.~Zhang, D.~del-Castillo-Negrete,
M.~Stoyanov, ``A Feynman--Kac based numerical method for the exit time
probability of a class of transport problems,'' \textit{J. Comput. Phys.}
\textbf{444}, 110564 (2021). doi:10.1016/j.jcp.2021.110564.

\end{thebibliography}

\end{document}
